\documentclass[11pt,fleqn]{report}
\usepackage{graphicx}
\usepackage{amssymb}
\usepackage{longtable}
\usepackage[T1]{fontenc}

\usepackage{cml}
\lstnewenvironment{codeXML}
    {\lstset{basicstyle=\small\ttfamily,
             language=XML,
             xleftmargin=0.0in,
             xrightmargin=0.0in,
             backgroundcolor=\color{Orchid!60}, % 60% orchid, 40% white
             escapeinside={\%*}{*)}, %  .tex file characters inside %*...*) are
                                     %  treated as Tex commands.
            }
    }
    { }

\newcommand{\ModelName}{Fault-Arch Legacy Interface}
\newcommand{\ModelAuthor}{Gary Turner}
\newcommand{\ModelAffiliation}{Odyssey Space Research}
\newcommand{\ModelDate}{Sept 2022}
\pagestyle{plain}
\usepackage[parfill]{parskip}
\usepackage[margin=1.5cm]{caption}
\usepackage{subfig}

\begin{document}
\titlePage
\clearpage % Reset page number
\pagenumbering{roman}
\tableofcontents % Print table of contents
\vfill
{\Large \textbf {Revision History}}
{\large
\begin{center}
\begin{tabular}{|l||l|l|l|}
\hline
 \textbf{Version} & \textbf{Date} & \textbf{Author} & \textbf{Purpose} \\ \hline
 \hline
1 & Sept 2022 & Gary Turner & Initial Version \\ \hline
\end{tabular}
\end{center}
}

\chapter{Introduction} % Make a "chapter" heading
\pagenumbering{arabic} % Start text with arabic 1
The Fault-Arch Legacy Interface model provides an interface to models using
the now-defunct \textit{FaultArch} model for their fault-injection
architecture. The \textit{FaultArch} model has been replaced with this
interface to the updated \textit{Fault Management} model.

For the most part, this new interface allows for the use of
the XML configuration files formerly used for a \textit{FaultArch}
implementation to still be used in the new \textit{Fault Management} model.

To the extent that we were able to identify features of the (undocumented)
\textit{FaultArch} model, those features have been tested in the new
\textit{Fault Management} model.

\chapter {Requirements}
The only requirement for this model is that it supports, with as little
intervention as possible, the use of legacy XML configuration files with the
new \textit{Fault Management} model.


\chapter {Model Specification}
The model provides a whole new definition for the legacy
\textit{sSensorFaults} class. The new implementation of
\textit{sSensorFaults}:
\begin{itemize}
\item inherits from the \textit{FaultManager} class in the \textit{Fault
Management} model
\item replaces the methods in the old \textit{sSensorFaults} with new
implementations that access methods in \textit{FaultManager}.
\item implements additional capabilities to translate syntax that has been
deprecated in \textit{FaultManager}.
\end{itemize}


\chapter {User's Guide}
This is intended to be a plug-and-play interface to support legacy XML files
configuring legacy fault-injection implementations.

There is no User's Guide for this model because the model is intentionally
not supposed to be ``used'' in any manual sense. To implement a
fault-injection system, use the \textit{FaultManager} model directly.


\chapter{Verification}
There will likely be some syntax or features that were used in the legacy
model that this interface does not support. The original model had no
documentation and it is impossible to guarantee that this model provides a
comprehensive coverage of the old model without knowledge of the syntax and
functional capabilities of the old model.

For model verification, we take a set of legacy XML files and use them to
configure the \textit{FaultManager} model via the new \textit{sSensorFaults}
implementation.

The legacy XML files that we tested work with the
\textit{FaultManager} model via this interface.

The \textit{ERROR\_code\_coverage} test uses \textit{xml\_files/error.xml},
which was NOT a legacy XML file; it was created to specifically test the
fault-detection capabilities of this interface. It will generate several
error messages.

\section{Code Coverage}
\begin{figure}[h!]
  \centering
  \includegraphics[scale=0.5]{graphics/code_coverage.png}
  \caption{Code Coverage}
  \label{fig:coverage}
\end{figure}

The verification comprises the following runs:
\begin{enumerate}
\item \textbf{RUN\_bias} applies 3 biases to 2 variables, starting at t=3,
      (+50 bias), t=5 (+100 bias), and t=8 (+150 bias).
      Bias faults accumulate when more than one is applied to the
      same variable, in this case the first and third are applied to the same
      variable to create a net bias of +200.
\begin{figure}[h!]
  \centering
  \includegraphics[scale=0.5]{graphics/bias.png}
  \caption{}
\end{figure}

\item \textbf{RUN\_linear} applies a linear walk-off starting at t=6.
\begin{figure}[h!]
  \centering
  \includegraphics[scale=0.5]{graphics/linear.png}
  \caption{}
\end{figure}

\item \textbf{RUN\_linear\_fault1} results in an error message and no fault
      applied.
\begin{figure}[h!]
  \centering
  \includegraphics[scale=0.5]{graphics/linear_fault.png}
  \caption{}
\end{figure}

\item \textbf{RUN\_linear\_fault2} results in an error message and no fault
      applied.
\begin{figure}[h!]
  \centering
  \includegraphics[scale=0.5]{graphics/linear_fault.png}
  \caption{}
\end{figure}

\newpage
\item \textbf{RUN\_linear\_fault3} tests the automatic assignment when
      \textit{b} (from $y = ax+b$) is omitted from the specification.
      This should result in the assumption of $b=0$ and data matching those
      from \textit{RUN\_linear}.
\begin{figure}[h!]
  \centering
  \includegraphics[scale=0.5]{graphics/linear.png}
  \caption{}
\end{figure}

\item \textbf{RUN\_linear\_repeated} tests the effect of having multiple
      triggers and multiple trigger groups. A linear walk-off is applied for
      $5 \leq t \leq 10$ (triggers $t \geq 5$ and $t \leq 10$
      are both satisfied in first trigger-group) and $t \geq 15$
      (solitary trigger satisfied in second trigger-group).
\begin{figure}[h!]
  \centering
  \includegraphics[scale=0.5]{graphics/linear_repeat.png}
  \caption{}
\end{figure}

\item \textbf{RUN\_negative\_periodic} applies an overwrite error when
      \textit{negative\_time} is less than -5 (i.e. $time \geq 5$). The
      overwrite is assigned for 1.5 seconds every 3 seconds thereafter.
\begin{figure}[h!]
  \centering
  \includegraphics[scale=0.5]{graphics/periodic.png}
  \caption{}
\end{figure}

\item \textbf{RUN\_overwrite} applies 3 different overwrite faults to 2
      variables. When two overwrite faults are applied to the same variable,
      only the second one ``sticks''. The three overwrites use the same
      specifications as were used for the earlier bias test, namely:
      for $t \geq 3$, a +50 overwrite, for $t \geq 5$, a +100 overwrite, and
      for $t \geq 8$, a +150 overwrite.
      The first and third overwrites are applied to the same variable, so the
      first is overwritten by the third for $t \geq 8$.
\begin{figure}[h!]
  \centering
  \includegraphics[scale=0.5]{graphics/overwrite.png}
  \caption{}
\end{figure}

\item \textbf{RUN\_periodic} matches \textit{RUN\_negative\_periodic} using
      a trigger monitoring positive time.
\begin{figure}[h!]
  \centering
  \includegraphics[scale=0.5]{graphics/periodic.png}
  \caption{}
\end{figure}
\newpage
\item \textbf{RUN\_persistence\_trigger} applies a set of four overwrite
      faults to a single variable. Normally, this would result in only the
      last applied overwrite remaining (as in the previous test case).
      In this case, however, the respective
      triggers are given progressively shorter fire-limits, so we see the
      fire limits restrict each trigger and the overwrite value return back
      through the set of overwrites as the later overwrites (with their shorter
      fire-times) expire in reverse order.
      The last fault remaining is the first fault applied, with a +100
      overwrite; when its trigger reaches its
      fire-limit, the variable is returned to its nominal value with no
      further faults to apply.
\begin{figure}[h!]
  \centering
  \includegraphics[scale=0.5]{graphics/persistence.png}
  \caption{}
\end{figure}

\item \textbf{RUN\_random} applies white-noise to the variable. The random
      distribution parameters are set by the recommended practice for the
      \textit{FaultManager} class.
\begin{figure}[h!]
  \centering
  \includegraphics[scale=0.5]{graphics/random.png}
  \caption{}
\end{figure}

\item \textbf{RUN\_random\_overwrite} generates a random value for an
      overwrite fault, then applies it consistently.
\begin{figure}[h!]
  \centering
  \includegraphics[scale=0.5]{graphics/random_ow.png}
  \caption{}
\end{figure}

\item \textbf{RUN\_random\_trigger} applies a fixed overwrite, starting at a
      random time. The random time aspect is confirmed by running several
      times with different seeds.
\begin{figure}[h!]
  \centering
  \includegraphics[scale=0.5]{graphics/random_trig.png}
  \caption{}
\end{figure}

\item \textbf{RUN\_random\_walk} applies a random walk fault.
\begin{figure}[h!]
  \centering
  \includegraphics[scale=0.5]{graphics/random_walk.png}
  \caption{}
\end{figure}

\item \textbf{RUN\_reset\_overwrite} tests the ability to change a parameter
      in the input file; the XML file sets the overwrite value to 50, and
      the input file sets it to 40. The results confirms the 40 is applied.
\begin{figure}[h!]
  \centering
  \includegraphics[scale=0.5]{graphics/reset_overwrite.png}
  \caption{}
\end{figure}

\newpage
\item \textbf{RUN\_reuse\_trigger} confirms that the same trigger can be
      used for multiple faults.
\begin{figure}[h!]
  \centering
  \includegraphics[scale=0.5]{graphics/reuse_trig.png}
  \caption{}
\end{figure}

\item \textbf{RUN\_scale} tests the scaling of some value. Three tests are
      included with three faults. The first test uses the variable
      presented in the middle column below; it scales the target variable
      by $\times 10$ and because this variable does not get reset
      externally, it keeps increasing by a factor of 10.

      The second test applies to the variable presented in the right-hand
      column below. It, too, applies a scaling factor of 10 but its target
      variable does get reset externally each cycle (to 8.675309), and then
      re-scaled (to 86.75309), resulting in the same faulted value each
      cycle.

      The third also applies a scaling of $\times 10$ to the variable in the
      right-hand column for $t \geq 8$. This scaling compounds with that
      from the second test, resulting in the scaled value (86.75309) being
      further scaled to 867.5309.
\begin{figure}[h!]
  \centering
  \includegraphics[scale=0.5]{graphics/scale.png}
  \caption{}
\end{figure}

\newpage
\item \textbf{RUN\_sinewave} applies a sinewave bias to the variable with a
      period of 8 seconds and amplitude 1.0.
\begin{figure}[h!]
  \centering
  \includegraphics[scale=0.5]{graphics/sine.png}
  \caption{}
\end{figure}

\newpage
\item \textbf{RUN\_sinewave\_var} applies three faults in sequence, with
      varying parameters:
  \begin{enumerate}
  \item frequency ($2 \leq t \leq 10$) -- oscillations become
        closer as the frequency increases;
  \item amplitude ($12 \leq t \leq 20$) -- oscillations maintain constant
        frequency but the amplitude grows linearly
  \item phase-offset ($t \geq 22$) -- osciallations are twice as fast as they
        would be otherwise as the phase offset increases at the same rate as
        the natural phase, effectively doubling the rate at which the net
        phase evolves.
  \end{enumerate}
\begin{figure}[hb]
  \centering
  \includegraphics[scale=0.5]{graphics/sine_var.png}
  \caption{}
\end{figure}

\clearpage
\item \textbf{RUN\_sinewave\_var\_repeated} applies the same three faults
      as the previous test case, then applies them again. The second set of
      applications matches the first. This confirms that the independent
      variable is being considered relative to its value at triggering;
      if its absolute value was being used, the second pass through each
      fault would be different from the first.
\begin{figure}[h!]
  \centering
  \includegraphics[scale=0.5]{graphics/sine_repeat.png}
  \caption{}
\end{figure}

\clearpage
\item \textbf{RUN\_squarewave} applies a bias with value from a squarewave
      function.
\begin{figure}[h!]
  \centering
  \includegraphics[scale=0.5]{graphics/square.png}
  \caption{}
\end{figure}

\item \textbf{RUN\_squarewave\_var} applies three faults in sequence, with
      varying frequency, amplitude, and phase-offset respectively. This is
      similar to the application of an evolving fault in
      \textit{RUN\_sinewave\_var} but for a squarewave rather than a
      sinewave. Note that the squarewave levels are not constant when the
      amplitude is varying.
\begin{figure}[h!]
  \centering
  \includegraphics[scale=0.5]{graphics/square_var.png}
  \caption{}
\end{figure}

\newpage
\item \textbf{RUN\_stale} applies stale-value overrides to
      \textit{counting\_int} and \textit{linear\_value} (the two values in
      the far right columns below).
\begin{figure}[h!]
  \centering
  \includegraphics[scale=0.5]{graphics/stale.png}
  \caption{}
\end{figure}

\item \textbf{RUN\_turnofffaults} repeats \textit{RUN\_sinewave\_var} but
      with all three faults switched off in the input file. None get
      applied.
\begin{figure}[h!]
  \centering
  \includegraphics[scale=0.5]{graphics/turnoff.png}
  \caption{}
\end{figure}

\newpage
\item \textbf{RUN\_turnofftriggers} repeats \textit{RUN\_sinewave\_var} but
      with all five triggers switched off in the input file. No faults get
      applied.
\begin{figure}[h!]
  \centering
  \includegraphics[scale=0.5]{graphics/turnoff.png}
  \caption{}
\end{figure}
\end{enumerate}



\end{document}
